% from main.tex
%%%%%%%%%%  scan idx 152  |  folio 139  %%%%%%%%%%
Autrement dit, pour que \(E\in\Ob\Dpx^b({}_A X)\) soit pseudo-cohérent il faut et il suffit que \(E\) soit parfait.

\emph{Preuve.} Soient \(E\in\Ob\Dpx^b({}_A X)_{\mathrm{coh}}\) et \(s\) un point de \(S_{/X}\). Il s'agit de montrer que \(E\) est parfait « au voisinage de \(s\) », i.e. qu'il existe un objet \(s\)-ponctué \(U\) de \(S\) au-dessus de \(X\) tel que \(E|U\) soit parfait. Or, le foncteur \(s^*\) étant exact, on a, d'après (2.16.1 a)), \(E_s\in\Ob\Dpx^b(A_s)_{\mathrm{coh}}\). Donc, en vertu de (5.9), \(E_s\) est parfait. L'assertion sera donc conséquence de (3.7.2) et du lemme suivant :

\textbf{Lemme 5.10.1.} \emph{Soit \(s\) un point de \(S\) et soit \(F\in\Ob\Dpx(A_s)_{\mathrm{parf}}\). Il existe un objet \(s\)-ponctué \(U\) de \(S\) et \(F'\in\Ob\Dpx({}_A U)_{\mathrm{parf}}\) tel que \(F'_s=F\).}

\emph{Preuve.} C'est évident si \(F\) est un \(A_s\)-module libre de type fini. Supposons maintenant que \(F\) soit projectif de type fini, donc image d'un projecteur \(p:L\longrightarrow L\) où \(L\) est libre de type fini. Il existe alors un « voisinage » \(U\) de \(s\), un \(A_U\)-Module libre \(L'\) et un homomorphisme \(p':L'\to L'\) tels que \(p'_s=p\). Comme \({p'_s}^2=p'_s\), on peut supposer, quitte à rétrécir \(U\), que \(p'^2=p'\). Donc \(F'=\Im(p')\) est facteur direct de \(L'\) et vérifie \(F'_s=F\). On généralise immédiatement au cas où \(F\) est strictement parfait, i.e. un complexe à degrés bornés composé de modules projectifs de type fini.

\textbf{Exemples 5.11.} Le corollaire (5.10) s'applique en particulier au cas où \((S,A)\) est défini par un espace annelé en anneaux locaux \emph{réguliers}, par exemple un schéma régulier, ou un espace analytique complexe lisse sur \(\mathbb C\).

